\chapter{Bài toán và công trình liên quan}
\label{ch:problem}

\section{Bối cảnh và bài toán thiết kế}
Phép nhân--tích lũy ($\sum_i A_i B_i$) là phép toán hạt nhân trong nhiều hệ thống
xử lý tín hiệu số (DSP), bộ lọc số, tính toán ma trận và các lớp suy luận
mạng nơ-ron. Vì phép toán được lặp lại với mật độ lớn, việc lựa chọn định dạng
biểu diễn số và cách tổ chức đường ống phần cứng ảnh hưởng trực tiếp đến tài nguyên
logic, tần số hoạt động, thời gian xử lý và sai số tích lũy của toàn chuỗi tính toán.

Trong thực tế công nghiệp, chuẩn dấu phẩy động IEEE-754 \cite{ieee754_2019} cung
cấp cơ sở tính toán phổ biến với quy tắc ngoại lệ chặt chẽ, nhưng số bit dành cho
trường mũ (\textit{exponent}) và phần định trị (\textit{fraction}) được ấn định cố
định cho từng định dạng (ví dụ 8 bit mũ và 23 bit fraction ở \texttt{binary32}).
Ngược lại, hệ thống số Posit do Gustafson và Yonemoto đề xuất \cite{gustafson2017posit}
và được chuẩn hóa năm 2022 tại \cite{positstandard2022} (với tham số chuẩn $es=2$)
sử dụng trường \textit{regime} có độ dài biến thiên trên cùng kích thước từ mã
32 bit. Cơ chế này cho phép phân bổ thêm bit fraction quanh miền $\pm 1$, đồng
thời cho phép mở rộng sang cấu hình tùy biến \texttt{Posit(32,3)} ($es=3$) theo
\cite{norris2021iterativeposit} để tăng mạnh dải động khi cần thiết. Tuy nhiên,
việc hiện thực trực tiếp bộ nhân Posit đầy đủ đòi hỏi bộ nhân cứng $27\times 27$ bit
kèm logic giải mã và tái mã hóa \textit{regime} phức tạp, gây tốn LUT và làm tăng
độ trễ tổ hợp trên FPGA.

Từ bài toán đánh đổi giữa chi phí phần cứng và độ chính xác số học, đề tài lựa
chọn hướng bộ nhân lặp dịch--cộng xấp xỉ có cấu hình giới hạn số vòng theo dữ liệu,
tích hợp trực tiếp vào hai kiến trúc nhân--cộng \texttt{MAC-NF} và \texttt{MAC-F}
cùng giao thức bắt tay \texttt{valid/ready}, đồng thời bổ sung kiến trúc tích lũy
trên thanh ghi mở rộng \texttt{MAC-Quire} để đối chứng sai số chuỗi dài.

\section{Mục tiêu và phạm vi}
Đồ án hiện thực và kiểm chứng chức năng lõi IP \texttt{Posit(32,3)} trên FPGA dòng
Artix-7 (\path{xc7a35tcpg236-1}), sử dụng độ rộng fraction nhân 12 bit, lưới tích
nội bộ \texttt{FULL\_TERM26/Q24} và quy tắc làm tròn chẵn gần nhất (RNE) tại khối
đóng gói Packer. Nội dung đánh giá phân tách rõ hai bài toán: giao dịch độc lập
$(A,B,C)$ và chuỗi tích lũy mà đầu vào $C$ được hồi tiếp từ kết quả $D$ của bước
trước. Trong ba biến thể kiến trúc, \textbf{MAC-F} (hợp nhất tích nội bộ sang bộ
cộng) là lõi thiết kế chính; \textbf{MAC-NF} (đóng gói tích trung gian về Posit
trước khi cộng) và \textbf{MAC-Quire} đóng vai trò đối chứng kiến trúc, tài nguyên,
tần số và độ chính xác theo hợp đồng thiết kế trình bày tại Chương~\ref{ch:design}.

\section{Công trình liên quan và phần việc của đồ án}
Hệ thống số Posit với quy tắc làm tròn RNE và bộ tích lũy chính xác \textit{Quire}
được Gustafson và Yonemoto đề xuất nhằm nâng cao độ chính xác cục bộ quanh $1$ và
triệt tiêu sai số làm tròn lặp giữa các phép MAC liên tiếp
\cite{gustafson2017posit,positstandard2022}. Tuy vậy, các phân tích so sánh giữa
Posit và chuẩn dấu phẩy động IEEE-754 trên phần cứng
\cite{dinechin2019posits,ciocirlan2021accuracy} cho thấy ưu thế lý thuyết của
Posit phải đánh đổi bằng chi phí giải mã và mã hóa \textit{regime} có độ dài biến
thiên. Trên nền tảng FPGA, các công trình tiêu biểu như bộ sinh lõi của Jaiswal
và So (\texttt{PACoGen}) \cite{jaiswal2018posit,jaiswal2019pacogen}, kiến trúc
tham số hóa của Chaurasiya và cộng sự \cite{chaurasiya2018parameterized}, thiết kế
của Podobas và Matsuoka \cite{podobas2018posit}, công cụ sinh lõi \texttt{FloPoCo}
của Murillo và cộng sự \cite{murillo2020flopoco,murillo2022energy}, bộ nhân tối ưu
công suất của Zhang và Ko \cite{zhang2020positmul} cùng khảo sát cho học sâu của
Lu và cộng sự \cite{lu2020evaluation} đều chỉ ra rằng khối dò bit dẫn đầu (LZD/LOD)
kết hợp hai bộ dịch đa tầng (\textit{barrel shifter}) và bộ nhân định trị $27\times 27$
bit là nguyên nhân chính làm tăng số LUT và độ trễ tổ hợp.

Để giảm mạnh chi phí bộ nhân phần cứng, nhiều công trình theo hướng tính toán xấp xỉ
và hỗn hợp độ chính xác cho mạng nơ-ron hoặc DSP đã được đề xuất, tiêu biểu là
\textit{Deep Positron} \cite{carmichael2019deep}, \textit{PositNN}
\cite{langroudi2019positnn}, bộ nhân xấp xỉ logarit \texttt{PLAM}
\cite{murillo2021plam}, bộ nhân xấp xỉ tiêu thụ thấp của Wan và Ko
\cite{wan2023lowcost} và các lõi xử lý tích hợp Quire như \texttt{PERCIVAL}
\cite{mallasen2022percival,mallasen2024bigposit}. Riêng với giải pháp nhân lặp
không dùng DSP, Norris và Kim \cite{norris2021iterativeposit} đề xuất bộ nhân
\texttt{Posit(32,3)} dịch--cộng cấu hình giới hạn vòng lặp, giữ 12 bit cao của
fraction với ba khối \texttt{OPS}, \texttt{SAC} và \texttt{SBM}. Phân định rõ ràng
giữa phần kế thừa và đóng góp mới: đồ án kế thừa nguyên lý giải mã--tái mã
hóa đường ống từ tài liệu trên và lõi nhân dịch--cộng xấp xỉ 12 bit của Norris--Kim,
từ đó phát triển các phần việc trọng tâm của đề tài gồm: (i) chuẩn hóa hợp đồng
bộ nhân lặp \texttt{FULL\_TERM26/Q24} (bảo toàn đủ 26 bit tích Q24 cùng hai cờ
cắt độc lập \texttt{input\_cut}/\texttt{approx\_cut}); (ii) thiết kế tích hợp trọn
vẹn ba kiến trúc \texttt{MAC-NF}, \texttt{MAC-F} (bỏ đóng gói trung gian) và
\texttt{MAC-Quire} với đường cộng 5 tầng co giãn trên cùng giao thức
\texttt{valid/ready}; (iii) phân tích nguyên nhân nghẽn định tuyến hậu route và
tối ưu tách tải fanout tầng S1 để đưa $F_{\max,\mathrm{R2R}}$ của \texttt{MAC-F}
lên $154{,}42$\,MHz (toàn thiết kế $129{,}0$\,MHz) với \texttt{0 DSP} trên Artix-7,
đặt trong đối chứng trực tiếp với IP \texttt{Floating-Point FP32} của AMD/Xilinx
\cite{amd_pg060_fp,amd_ug906_timing}; và (iv) xây dựng mô hình số nguyên lớn đối
chiếu độc lập với \texttt{SoftPosit} kèm đánh giá ứng dụng thực tế và độ phủ cuối cùng.

\section{Hai bài toán sử dụng cùng giao diện MAC}

Trong bài toán thứ nhất, mỗi bộ ba $(A,B,C)$ là một giao dịch độc lập.
Bộ điều khiển bên ngoài có thể cấp $C$ tùy theo dữ liệu ứng dụng; $C$ không
mang nghĩa trạng thái bên trong RTL của MAC-NF/F. Ở cấu hình mặc định, hai bộ
nhân nhận cùng mã $A/B$, giữ cùng phần fraction và phát cùng tích nội bộ. Do
đó có thể sử dụng cặp kết quả $(D_{\mathrm{NF}}, D_{\mathrm{F}})$ trên cùng
dữ liệu đầu vào để quan sát trực tiếp tác động của vị trí làm tròn trung gian.

Trong bài toán thứ hai, trạng thái $S_i$ của một tích vô hướng được dùng làm
toán hạng cộng cho bước tiếp theo. Bộ điều khiển ngoài của MAC-NF hoặc MAC-F
đưa $C = S_{i-1}$ và chỉ cập nhật $S_i \gets D_i$ khi giao dịch trước hoàn
tất (\textit{retire}). Nếu hai lõi cho kết quả khác nhau tại một bước, các giá
trị $C$ hồi tiếp về sau cũng phân kỳ theo. Phép đánh giá lúc này là đối chiếu
cả chuỗi trên cùng dãy toán hạng đầu vào $(A_i, B_i)$, cùng điều kiện bắt tay
và trạng thái khởi tạo ban đầu. Đối với \texttt{MAC-Quire}, trạng thái tích lũy
được giữ nội bộ dưới dạng thanh ghi 972 bit với trình tự một lệnh \texttt{CLEAR},
chuỗi lệnh \texttt{MAC} và một lệnh \texttt{READ} cuối chuỗi (chi phí chu kỳ của
\texttt{CLEAR} và \texttt{READ} được tính gộp vào tổng chu kỳ công bố).

\blockfigure{Hai mức đo trên cùng giao diện số học. Mũi tên đứt là đường hồi tiếp bên ngoài của MAC-NF và MAC-F.}{fig:measurement-levels}{
  \node[block] (inputs) {$A_i,B_i,C_i$\\giao dịch độc lập};
  \node[block,right=7mm of inputs] (dut) {MAC\\NF hoặc F};
  \node[block,right=7mm of dut] (output) {$D_i$, cờ\\khi retire};
  \node[store,below=7mm of dut] (state) {State driver\\$S_i\gets D_i$};
  \draw[data] (inputs)--(dut);
  \draw[data] (dut)--(output);
  \draw[control] (output.south)|-(state.east);
  \draw[control] (state.west)-|(inputs.south) node[pos=.22,below] {\scriptsize chế độ chuỗi};
}

\section{Câu hỏi đánh giá và bằng chứng}

\begin{engineeringtable}{Câu hỏi đánh giá và phép đối chiếu tương ứng.}{tab:research-questions}
\begin{tabularx}{\linewidth}{@{}>{\raggedright\arraybackslash}p{42mm}Y Y@{}}
\toprule
Câu hỏi & Phép đo & Phạm vi suy luận \\
\midrule
Vòng lặp tạo đúng tích quy định? & So $M$, dấu, $sf$, số vòng và cờ cắt với golden ở \texttt{core\_done}. & Xác nhận tính đúng số học trước chuẩn hóa và RNE. \\
Sai số phát sinh từ đâu? & So cùng đầu vào với tích đủ fraction và tích của hai significand đã cắt 12 bit. & Tách riêng sai số cắt đầu vào và sai số giới hạn vòng. \\
Nên chọn NF hay F? & Đo ghép cặp $(A,B,C)$, tổng hợp và định tuyến trên cùng part FPGA và ràng buộc clock. & Lựa chọn kiến trúc ở cấp IP dưới các điều kiện đã ghi. \\
Làm tròn giữa các bước ảnh hưởng thế nào? & So cùng dãy tích, tổng số chu kỳ và sai số cuối của NF, F và quire. & Đánh giá hiệu năng cho các nhóm chuỗi đã mô phỏng. \\
\bottomrule
\end{tabularx}
\end{engineeringtable}

Tính đúng thiết kế (\textit{functional correctness}) và độ chính xác số học
(\textit{arithmetic accuracy}) được đánh giá trên hai hệ mốc độc lập. Một lõi
MAC có thể khớp 100\% về từng bit với mô hình tham chiếu xấp xỉ nhưng kết quả
vẫn lệch một khoảng đo được so với giá trị thực chính xác $AB+C$. Tương tự, kết
quả tổng hợp và định tuyến xác định tần số và tài nguyên phần cứng, không thay
thế cho kiểm chứng chức năng. Chương~\ref{ch:evaluation} trình bày phương pháp
và kết quả kiểm chứng chức năng; Chương~\ref{ch:measurements} trình bày các đại
lượng định lượng về sai số, chu kỳ và triển khai FPGA.

\section{Đối chiếu hướng thiết kế kế thừa}

Nghiên cứu của Norris--Kim sử dụng độ rộng fraction cố định ở ngõ vào bộ nhân và
cho phép chọn giới hạn số vòng lặp khi chạy. Đồ án kế thừa ý tưởng liệt kê các số
hạng dịch từ bit có trọng số cao xuống thấp, đồng thời cụ thể hóa đầy đủ hợp đồng
kỹ thuật để tích hợp vào MAC: chốt độ rộng tích nội bộ Q24 (26 bit không dấu),
quy định chính sách chọn toán hạng \texttt{OPS} theo số bit~1 ít hơn, bổ sung cờ
báo mất bit ở hai nguồn (cắt đầu vào và giới hạn vòng) và xác định ranh giới làm
tròn RNE. Ngoài ra, thay vì chỉ báo tỷ lệ đạt ngưỡng sai số trên kết quả đã đóng
gói khi chuyển đổi từ chuẩn dấu phẩy động IEEE-754 như trong \cite{norris2021iterativeposit},
đồ án đo trực tiếp trên bộ mã \texttt{Posit(32,3)} cả trước và sau Packer để
phân tách tường minh hai nguồn sai số.

\section{Phạm vi thiết kế và bố cục báo cáo}

Toàn bộ báo cáo sử dụng cấu hình mở rộng \texttt{Posit(32,3)} ($es=3$) thống nhất trong lý thuyết,
ví dụ tính tay, mô hình tham chiếu (\textit{golden model}) và RTL. Cấu hình đánh
giá mặc định đặt \texttt{OPS=1} và \texttt{n=12}; các giá trị cấu hình còn lại
(\texttt{OPS=0}, $n\in[0,12]$) được kiểm chứng đầy đủ trong hồi quy hỗn hợp và
khảo sát quét sai số bộ nhân. Các sơ đồ datapath trong Chương~\ref{ch:design} phản
ánh trực tiếp mã RTL cuối cùng đã vượt qua kiểm chứng mô phỏng và đóng định thời
sau \textit{Place \& Route} OOC ở 100\,MHz.

Báo cáo được tổ chức thành 6 chương: Chương~\ref{ch:problem} đặt bài toán và
liên hệ các công trình liên quan; Chương~\ref{ch:arithmetic} trình bày cơ sở
toán học của \texttt{Posit(32,3)}, lý do lựa chọn thay cho IEEE-754 cùng hệ thống
ví dụ giải mã, nhân, cộng/trừ từng bước; Chương~\ref{ch:design} mô tả chi tiết
vi kiến trúc RTL của lõi chính \texttt{MAC-F} và hai cấu hình đối chứng;
Chương~\ref{ch:evaluation} trình bày phương pháp kiểm chứng và kết quả mô phỏng
chức năng; Chương~\ref{ch:measurements} đánh giá định lượng sai số, chu kỳ chuỗi
tích lũy và kết quả triển khai FPGA kèm tối ưu hóa đường găng sau định tuyến; và
Chương~\ref{ch:conclusion} tổng kết kết quả cùng hướng phát triển.

\section{Đóng góp của đề tài}

Đóng góp thứ nhất là hoàn thiện hợp đồng định lượng cho bộ nhân lặp xấp xỉ
\texttt{FULL\_TERM26/Q24}. Độ rộng 12 bit fraction đầu vào, thanh ghi tích lũy
26 bit lưới Q24 và hai cờ mất dữ liệu tách biệt cho phép kiểm tra trực tiếp
từng kết quả tại ranh giới khối nhân và tái lập hoàn toàn kết quả đo theo $n$.

Đóng góp thứ hai là thiết kế và tích hợp cùng khối nhân lặp vào hai kiến trúc
\texttt{MAC-NF} và \texttt{MAC-F} trên cùng chuẩn giao thức \texttt{valid/ready},
kèm theo phiên bản đối chứng \texttt{MAC-Quire} sử dụng ngân hàng tích lũy 972 bit
theo từng từ 36 bit và giao thức lệnh \texttt{CLEAR/MAC/READ}.

Đóng góp thứ ba là quá trình phân tích sâu nguyên nhân sụt giảm định thời giữa pha
tổng hợp và định tuyến trên Artix-7, từ đó đề xuất giải pháp tách tải điều khiển
fanout cao (\texttt{choose\_a}) tại tầng S1 mà không làm thay đổi datapath, số FF
hay độ trễ chu kỳ của kiến trúc trọng tâm \texttt{MAC-F}.

\begin{engineeringtable}{Phạm vi suy luận và điều kiện áp dụng của các nhóm kết quả.}{tab:claim-levels}
\begin{tabularx}{\linewidth}{@{}>{\raggedright\arraybackslash}p{34mm}Y Y@{}}
\toprule
Nhóm kết quả & Kết luận rút ra trực tiếp & Điều kiện khi mở rộng hệ thống \\
\midrule
Kiểm chứng chức năng & Payload 32 bit và 5 cờ trạng thái khớp 100\% mô hình số nguyên lớn trên toàn bộ tập kiểm thử. & Mở rộng ma trận kiểm chứng quy mô lớn khi tích hợp vào SoC hoàn chỉnh. \\
Đánh giá sai số & Phân bố sai số cắt đầu vào, sai số vòng lặp và sai số ghép cặp NF/F trên tập lấy mẫu đều, phân tầng và triệt tiêu. & Hiệu chỉnh thống kê nếu áp dụng cho phân bố trọng số của một mạng nơ-ron cụ thể. \\
Triển khai FPGA (OOC) & Tài nguyên, độ trễ F2F, WNS/WHS, $F_{\max,\mathrm{R2R}}$ và $F_{\max}$ toàn TK dưới ràng buộc OOC 100\,MHz. & Tái kiểm tra timing tổng thể sau khi gắn chân I/O vật lý và đăng ký cổng \texttt{in\_ready}. \\
\bottomrule
\end{tabularx}
\end{engineeringtable}

